\subsection{Asynchronous Messaging (RabbitMQ)}

The platform uses a single topic exchange, \texttt{clinical.events}, with routing keys of the form \texttt{<entity>.<action>}: \texttt{patient.created}, \texttt{provider.created}, \texttt{encounter.created}, \texttt{notification.send}, \texttt{password.reset\_requested}, \texttt{cluster.run\_requested}. \texttt{auth} subscribes to \texttt{patient.created} / \texttt{provider.created} (account provisioning) and \texttt{password.reset\_requested}; \texttt{mail-service} subscribes to \texttt{notification.send}; \texttt{ml-batch} subscribes to \texttt{cluster.run\_requested}. The choice of \emph{topic} over \emph{fanout} is deliberate: new event types are introduced without touching consumers that should not receive them. Producers use a uniform \texttt{publish\_event(routing\_key, payload)} helper that serialises to JSON, sets persistent delivery and reuses one channel per worker.
